% GENERATED FILE. Edit canonical lessons, then run npm run generate:books.
% canonical-source-hash: fnv1a64:eec76bef7a43947c
% canonical-lessons: SA-C548-visvasayati, SA-C548-durdasa, SA-C548-santvana, SA-C548-akulah, SA-C548-saranam, SA-R548-review-542-548

\chapter{Fears and Reassurance}
\label{ch:sa-fears-and-reassurance}
\hlchaptermodality{548}

\begin{chapteropening}
\textbf{By the end of this chapter:} \emph{I can comfort someone who is afraid of what might happen.}

Complete the last lesson of chapter 548: I can comfort someone who is afraid of what might happen.
\end{chapteropening}

\section[\texorpdfstring{\sk{विश्वासयति} (viśvāsayati)}{विश्वासयति (viśvāsayati)}]{\sk{विश्वासयति} (viśvāsayati) — inspires confidence}

\label{lesson:SA-C548-visvasayati}



\begin{quote}
Before the new one: say the Sanskrit for a climb, then the Sanskrit for springs up.
\end{quote}

\subsection*{You'll want to know: \sk{विश्वासयति}}
\textbf{\sk{विश्वासयति}} — \emph{viśvāsayati} — \textquotedblleft{}inspires confidence\textquotedblright{}.

\begin{quote}
\sk{यदि} \sk{बालः} \sk{भीतः} \sk{अस्ति}, \sk{तर्हि} \sk{माता} \sk{तं} \sk{विश्वासयति।}
\end{quote}

\begin{quote}
\emph{yadi bālaḥ bhītaḥ asti, tarhi mātā taṃ viśvāsayati.}
\end{quote}

\begin{quote}
If the child is afraid, my mother gives him confidence.
\end{quote}

Past (she …): \sk{सा} \sk{बालं} \sk{विश्वासितवती}, \emph{sā bālaṃ viśvāsitavatī}.

\begin{cousinweb}
\textbf{\sk{विश्वासयति}} (\emph{viśvāsayati}) is built on \emph{vi} and the causative of the root √śvas, \textquotedblleft{}to breathe\textquotedblright{}, the root of \emph{viśvasiti}, \textquotedblleft{}trusts\textquotedblright{}; MW gives \emph{viśvāsayati} as to cause to trust, to inspire with confidence, to console (MW).
\end{cousinweb}

\begin{grammarlens}[title={What you've built}]
One more word for saying what happens if.
\end{grammarlens}

\subsection*{Your turn}
\begin{itemize}
\raggedright
  \item \emph{Say it:} \emph{viśvāsayati}
  \item \emph{Say it:} \emph{viśvāsayati}, once more
  \item \emph{Say it:} \emph{yadi bālaḥ bhītaḥ asti, tarhi mātā taṃ viśvāsayati.} — If the child is afraid, my mother gives him confidence.
  \item \emph{Recall:} say the Sanskrit for a climb, then the Sanskrit for springs up, then say \emph{viśvāsayati} again
\end{itemize}

\begin{tcolorbox}[breakable,colback=teal!4,colframe=teal!35!black,title={Before you move on}]
What does \sk{विश्वासयति} mean? (\textquotedblleft{}Inspires confidence\textquotedblright{}.) Say it once more.
\end{tcolorbox}
\section[\texorpdfstring{\sk{दुर्दशा} (durdaśā)}{दुर्दशा (durdaśā)}]{\sk{दुर्दशा} (durdaśā) — a sorry plight}

\label{lesson:SA-C548-durdasa}



\begin{quote}
Before the new one: say the Sanskrit for springs up, then the Sanskrit for inspires confidence.
\end{quote}

\subsection*{You'll want to know: \sk{दुर्दशा}}
\textbf{\sk{दुर्दशा}} — \emph{durdaśā} — \textquotedblleft{}a sorry plight\textquotedblright{}.

\begin{quote}
\sk{यदि} \sk{दुर्दशा} \sk{आगच्छति}, \sk{तर्हि} \sk{मित्रं} \sk{साहाय्यं} \sk{करोति।}
\end{quote}

\begin{quote}
\emph{yadi durdaśā āgacchati, tarhi mitraṃ sāhāyyaṃ karoti.}
\end{quote}

\begin{quote}
If hard times come, a friend helps.
\end{quote}

A feminine noun.

\begin{cousinweb}
\textbf{\sk{दुर्दशा}} (\emph{durdaśā}) is built on \emph{dur}, \textquotedblleft{}bad\textquotedblright{}, and \emph{daśā}, \textquotedblleft{}a state, a condition\textquotedblright{}; MW gives \emph{durdaśā} as a bad situation, misfortune (MW).
\end{cousinweb}

\begin{grammarlens}[title={What you've built}]
One more word for saying what happens if.
\end{grammarlens}

\subsection*{Your turn}
\begin{itemize}
\raggedright
  \item \emph{Say it:} \emph{durdaśā}
  \item \emph{Say it:} \emph{durdaśā}, once more
  \item \emph{Say it:} \emph{yadi durdaśā āgacchati, tarhi mitraṃ sāhāyyaṃ karoti.} — If hard times come, a friend helps.
  \item \emph{Recall:} say the Sanskrit for springs up, then the Sanskrit for inspires confidence, then say \emph{durdaśā} again
\end{itemize}

\begin{tcolorbox}[breakable,colback=teal!4,colframe=teal!35!black,title={Before you move on}]
What does \sk{दुर्दशा} mean? (\textquotedblleft{}A sorry plight\textquotedblright{}.) Say it once more.
\end{tcolorbox}
\section[\texorpdfstring{\sk{सान्त्वना} (sāntvanā)}{सान्त्वना (sāntvanā)}]{\sk{सान्त्वना} (sāntvanā) — consolation}

\label{lesson:SA-C548-santvana}



\begin{quote}
Before the new one: say the Sanskrit for inspires confidence, then the Sanskrit for a sorry plight.
\end{quote}

\subsection*{You'll want to know: \sk{सान्त्वना}}
\textbf{\sk{सान्त्वना}} — \emph{sāntvanā} — \textquotedblleft{}consolation\textquotedblright{}.

\begin{quote}
\sk{यदि} \sk{भ्राता} \sk{दुःखी} \sk{अस्ति}, \sk{तर्हि} \sk{तस्मै} \sk{सान्त्वनां} \sk{ददातु।}
\end{quote}

\begin{quote}
\emph{yadi bhrātā duḥkhī asti, tarhi tasmai sāntvanāṃ dadātu.}
\end{quote}

\begin{quote}
If your brother is unhappy, offer him comfort.
\end{quote}

A feminine noun.

\begin{cousinweb}
\textbf{\sk{सान्त्वना}} (\emph{sāntvanā}) comes from the root √sāntv, \textquotedblleft{}to console, to conciliate\textquotedblright{}; MW gives \emph{sāntvana}, beside \emph{sāntvanā}, as soothing with kind words, consolation (MW).
\end{cousinweb}

\begin{grammarlens}[title={What you've built}]
One more word for saying what happens if.
\end{grammarlens}

\subsection*{Your turn}
\begin{itemize}
\raggedright
  \item \emph{Say it:} \emph{sāntvanā}
  \item \emph{Say it:} \emph{sāntvanā}, once more
  \item \emph{Say it:} \emph{yadi bhrātā duḥkhī asti, tarhi tasmai sāntvanāṃ dadātu.} — If your brother is unhappy, offer him comfort.
  \item \emph{Recall:} say the Sanskrit for inspires confidence, then the Sanskrit for a sorry plight, then say \emph{sāntvanā} again
\end{itemize}

\begin{tcolorbox}[breakable,colback=teal!4,colframe=teal!35!black,title={Before you move on}]
What does \sk{सान्त्वना} mean? (\textquotedblleft{}Consolation\textquotedblright{}.) Say it once more.
\end{tcolorbox}
\section[\texorpdfstring{\sk{आकुलः} (ākulaḥ)}{आकुलः (ākulaḥ)}]{\sk{आकुलः} (ākulaḥ) — flustered; agitated}

\label{lesson:SA-C548-akulah}



\begin{quote}
Before the new one: say the Sanskrit for a sorry plight, then the Sanskrit for consolation.
\end{quote}

\subsection*{You'll want to know: \sk{आकुलः}}
\textbf{\sk{आकुलः}} — \emph{ākulaḥ} — \textquotedblleft{}flustered; agitated\textquotedblright{}.

\begin{quote}
\sk{यदि} \sk{भवान्} \sk{आकुलः} \sk{अस्ति}, \sk{तर्हि} \sk{माम्} \sk{आह्वयतु।}
\end{quote}

\begin{quote}
\emph{yadi bhavān ākulaḥ asti, tarhi mām āhvayatu.}
\end{quote}

\begin{quote}
If you are flustered, please call me.
\end{quote}

\begin{cousinweb}
\textbf{\sk{आकुलः}} (\emph{ākulaḥ}) is a word MW traces, with a query, to \emph{ā} and the root √kṝ, \textquotedblleft{}to scatter\textquotedblright{}; MW gives \emph{ākula} as confounded, confused, agitated, flurried, so its origin is uncertain (MW).
\end{cousinweb}

\begin{grammarlens}[title={What you've built}]
One more word for saying what happens if.
\end{grammarlens}

\subsection*{Your turn}
\begin{itemize}
\raggedright
  \item \emph{Say it:} \emph{ākulaḥ}
  \item \emph{Say it:} \emph{ākulaḥ}, once more
  \item \emph{Say it:} \emph{yadi bhavān ākulaḥ asti, tarhi mām āhvayatu.} — If you are flustered, please call me.
  \item \emph{Recall:} say the Sanskrit for a sorry plight, then the Sanskrit for consolation, then say \emph{ākulaḥ} again
\end{itemize}

\begin{tcolorbox}[breakable,colback=teal!4,colframe=teal!35!black,title={Before you move on}]
What does \sk{आकुलः} mean? (\textquotedblleft{}Flustered; agitated\textquotedblright{}.) Say it once more.
\end{tcolorbox}
\section[\texorpdfstring{\sk{शरणम्} (śaraṇam)}{शरणम् (śaraṇam)}]{\sk{शरणम्} (śaraṇam) — a refuge; shelter}

\label{lesson:SA-C548-saranam}



\begin{quote}
Before the new one: say the Sanskrit for consolation, then the Sanskrit for flustered.
\end{quote}

\subsection*{You'll want to know: \sk{शरणम्}}
\textbf{\sk{शरणम्}} — \emph{śaraṇam} — \textquotedblleft{}a refuge; shelter\textquotedblright{}.

\begin{quote}
\sk{यदि} \sk{वात्या} \sk{आगच्छति}, \sk{तर्हि} \sk{वयं} \sk{शरणं} \sk{गच्छामः।}
\end{quote}

\begin{quote}
\emph{yadi vātyā āgacchati, tarhi vayaṃ śaraṇaṃ gacchāmaḥ.}
\end{quote}

\begin{quote}
If a storm comes, we go to shelter.
\end{quote}

A neuter noun.

\begin{cousinweb}
\textbf{\sk{शरणम्}} (\emph{śaraṇam}) comes from the root √śri, \textquotedblleft{}to lean on, to resort to\textquotedblright{}, as MW derives it (through a form √śṛ); MW gives \emph{śaraṇa} as protecting, and a shelter, a place of refuge or rest, a house (MW).
\end{cousinweb}

\begin{grammarlens}[title={What you've built}]
One more word for saying what happens if.
\end{grammarlens}

\subsection*{Your turn}
\begin{itemize}
\raggedright
  \item \emph{Say it:} \emph{śaraṇam}
  \item \emph{Say it:} \emph{śaraṇam}, once more
  \item \emph{Say it:} \emph{yadi vātyā āgacchati, tarhi vayaṃ śaraṇaṃ gacchāmaḥ.} — If a storm comes, we go to shelter.
  \item \emph{Recall:} say the Sanskrit for consolation, then the Sanskrit for flustered, then say \emph{śaraṇam} again
\end{itemize}

\begin{tcolorbox}[breakable,colback=teal!4,colframe=teal!35!black,title={Before you move on}]
What does \sk{शरणम्} mean? (\textquotedblleft{}A refuge; shelter\textquotedblright{}.) Say it once more.
\end{tcolorbox}
\section[(dialogue)]{Looking back: from Making Up One's Mind to Fears and Reassurance}

\label{lesson:SA-R548-review-542-548}



\begin{quote}
Say the Sanskrit for acceptance and makes sure.
\end{quote}

\subsection*{Your turn}
Say these aloud:

\begin{itemize}
\raggedright
  \item acceptance, makes sure, excellent, pleasing, neutral
  \item free will, helpless, obstinacy, insistence, by oneself
  \item each other, at the same time, together, separately, at present
  \item a sequence, in due order, except, at first, reviews
  \item satisfied, completed, arises, an attainment, a consequence
  \item a rise, a decline, a fall, a climb, springs up
  \item inspires confidence, a sorry plight, consolation, flustered, a refuge
\end{itemize}

\begin{tcolorbox}[breakable,colback=teal!4,colframe=teal!35!black,title={Before you move on}]
Acceptance? (\textbf{\sk{स्वीकारः}}, \emph{svīkāraḥ}.) Makes sure? (\textbf{\sk{अवधारयति}}, \emph{avadhārayati}.) Excellent? (\textbf{\sk{उत्कृष्टः}}, \emph{utkṛṣṭaḥ}.) Pleasing? (\textbf{\sk{रुचिरः}}, \emph{ruciraḥ}.) Neutral? (\textbf{\sk{तटस्थः}}, \emph{taṭasthaḥ}.) Free will? (\textbf{\sk{स्वेच्छा}}, \emph{svecchā}.) Helpless? (\textbf{\sk{विवशः}}, \emph{vivaśaḥ}.) Obstinacy? (\textbf{\sk{हठः}}, \emph{haṭhaḥ}.) Insistence? (\textbf{\sk{आग्रहः}}, \emph{āgrahaḥ}.) By oneself? (\textbf{\sk{स्वयम्}}, \emph{svayam}.) Each other? (\textbf{\sk{परस्परम्}}, \emph{parasparam}.) At the same time? (\textbf{\sk{युगपत्}}, \emph{yugapat}.) Together? (\textbf{\sk{मिथः}}, \emph{mithaḥ}.) Separately? (\textbf{\sk{पृथक्}}, \emph{pṛthak}.) At present? (\textbf{\sk{सम्प्रति}}, \emph{samprati}.) A sequence? (\textbf{\sk{अनुक्रमः}}, \emph{anukramaḥ}.) In due order? (\textbf{\sk{यथाक्रमम्}}, \emph{yathākramam}.) Except? (\textbf{\sk{ऋते}}, \emph{ṛte}.) At first? (\textbf{\sk{आदौ}}, \emph{ādau}.) Reviews? (\textbf{\sk{आलोचयति}}, \emph{ālocayati}.) Satisfied? (\textbf{\sk{सन्तुष्टः}}, \emph{santuṣṭaḥ}.) Completed? (\textbf{\sk{निष्पन्नः}}, \emph{niṣpannaḥ}.) Arises? (\textbf{\sk{उद्भवति}}, \emph{udbhavati}.) An attainment? (\textbf{\sk{प्राप्तिः}}, \emph{prāptiḥ}.) A consequence? (\textbf{\sk{विपाकः}}, \emph{vipākaḥ}.) A rise? (\textbf{\sk{उत्थानम्}}, \emph{utthānam}.) A decline? (\textbf{\sk{अवनतिः}}, \emph{avanatiḥ}.) A fall? (\textbf{\sk{पतनम्}}, \emph{patanam}.) A climb? (\textbf{\sk{आरोहः}}, \emph{ārohaḥ}.) Springs up? (\textbf{\sk{उच्छलति}}, \emph{ucchalati}.) Inspires confidence? (\textbf{\sk{विश्वासयति}}, \emph{viśvāsayati}.) A sorry plight? (\textbf{\sk{दुर्दशा}}, \emph{durdaśā}.) Consolation? (\textbf{\sk{सान्त्वना}}, \emph{sāntvanā}.) Flustered? (\textbf{\sk{आकुलः}}, \emph{ākulaḥ}.) A refuge? (\textbf{\sk{शरणम्}}, \emph{śaraṇam}.)
\end{tcolorbox}
